% Titel etc
\title{Ein FAIRer Comic über Forschungsdateninfrastruktur}
\subtitle{Teil 1: Von Karl dem Großen zu ing.grid}
\shorttitle{Ein FAIRer Comic}

% Seite 1
\FDISetText{1}{1}{1}{Von Karl dem Großen bis zu Karl dem Kleinen.\\ Von der Karolingischen Minuskel über das Mittelreich bis zu \hspace*{0.5cm} NFDI4ING.* Forschungsdateninfrastruktur ist und bleibt eine spannende Sache.}
\FDISetText{1}{1}{2}{Westfrankenreich}
\FDISetText{1}{1}{3}{Mittelreich}
\FDISetText{1}{1}{4}{Ostfrankenreich}
\FDISetText{1}{1}{5}{Nordsee}
\FDISetText{1}{1}{6}{Baltisches\\Meer}
\FDISetText{1}{1}{7}{Kgr.\\England}
\FDISetText{1}{1}{8}{Irland}
\FDISetText{1}{1}{9}{Atlantischer\\Ozean}
\FDISetText{1}{1}{10}{Kgr. Asturien}
\FDISetText{1}{1}{11}{Bruessel}
\FDISetText{1}{1}{12}{Aachen}
\FDISetText{1}{1}{13}{Luxem\FDIStyledText{8pt}{2pt}{-}\\[-2pt]burg}
\FDISetText{1}{1}{14}{Verdun}
\FDISetText{1}{1}{15}{Darmstadt}
\FDISetText{1}{1}{16}{Strasburg}
\FDISetText{1}{1}{17}{Rom}
\FDISetText{1}{1}{18}{Emirat von Cordova}
\FDISetText{1}{1}{19}{Reich von Fez}
\FDISetText{1}{1}{20}{Mittelländisches Meer}
\FDISetText{1}{1}{21}{Benevent}
\FDISetText{1}{1}{22}{Serbien}
\FDISetText{1}{1}{23}{Slavische Völker}
\FDISetText{1}{1}{24}{* NFDI4ING ist das Konsortium für die Ingenieurwissenschaften in der Nationalen Forschungsdateninfrastruktur (NFDI)}
\FDISetText{1}{1}{25}{Limes Sorabicus}


% Seite 2
\FDISetText{2}{1}{1}{Im 8/9. Jahrhundert begann Karl der Große mit seiner zentralisierten Schriftkultur eine nachhaltige Wissensinfrastruktur aufzubauen. Er stützte sich dabei auf ein Netzwerk von Skriptoren in Klöstern und an seinem Hof in Aachen.}
\FDISetText{2}{2}{1}{Eine andere\\Sprache zu beherrschen\\ ist wie eine zweite Seele\\ zu besitzen.}
\FDISetText{2}{3}{1}{Karl selbst konnte nicht schreiben. Beim Unterschreiben von Dokumenten fügte er dem Karolus-Monogramm lediglich den sogenannten Vollziehungsstrich hinzu.}
\FDISetText{2}{4}{1}{Karl förderte aber Bildung durch das Diktieren, Kopieren und die Verbreitung von Dokumenten durch Gelehrte.}
\FDISetText{2}{4}{2}{Hefte raus-\\Klassenarbeit!}
\FDISetText{2}{5}{1}{Um die Lesbarkeit und Einheitlichkeit von Dokumenten in seinem riesigen Reich zu gewährleisten, ließ er die \enquote{Karolingigsche Minuskel} entwickeln.}
\FDISetText{2}{5}{2}{Nur kleine Buchstaben (Minuskel)}
\FDISetText{2}{5}{3}{Gleichmäßiges, offenes, ruhiges Schriftbild}
\FDISetText{2}{5}{4}{Das \enquote{i} wurde noch ohne Punkt geschrieben}
\FDISetText{2}{5}{5}{Vierlinien-System für bessere Lesbarkeit}
\FDISetText{2}{5}{6}{karl ist doof}
\FDISetText{2}{6}{1}{Das \enquote{s} wurde bis\\ins 12. Jahhundert\\als \enquote{langes s}\\geschrieben}
\FDISetText{2}{6}{2}{Der Buchstabe\\\enquote{j} existierte im 8.\\und 9. Jahrhundert\\noch nicht}
\FDISetText{2}{6}{3}{Unverbundene Buchstaben für bessere Lesbarkeit}
\FDISetText{2}{6}{4}{\FDIUseText{2}{5}{6}}
\FDISetText{2}{6}{5}{stimmt ia\\gar nicht}

% Seite 3
\FDISetText{3}{1}{1}{Karls Söhne und Enkel setzten diese analoge Forschungsinfrastruktur auf Basis von Pergament und Tinte fort.}
\FDISetText{3}{1}{2}{Karls Sippe}
\FDISetText{3}{1}{3}{(Eine Auswahl)}
\FDISetText{3}{1}{4}{Karl der Große}
\FDISetText{3}{1}{5}{Pippin\\[-2.5pt]der\\[-2.5pt]Bucklige}
\FDISetText{3}{1}{6}{Karl\\[-2.5pt]der\\[-2.5pt]Jüngere}
\FDISetText{3}{1}{7}{Karlmann}
\FDISetText{3}{1}{8}{Ludwig\\[-2.5pt]der\\[-2.5pt]Fromme}
\FDISetText{3}{1}{9}{Lothar I.}
\FDISetText{3}{1}{10}{Pippin}
\FDISetText{3}{1}{11}{Ludwig II.\\[-2.5pt]der\\[-2.5pt]Deutsche}
\FDISetText{3}{1}{12}{Karl II.\\[-2.5pt]der\\[-2.5pt]Kahle}
\FDISetText{3}{2}{1}{Einer seiner Enkel, \FDIUseText{3}{1}{9}, war im\\Jahr 843 maßgeblich an dem Vertrag\\von Verdun beteiligt.}
\FDISetText{3}{3}{1}{Dieser Vertrag gilt als Geburtsstunde der Vorläufer von Frankreich und Deutschland (\FDIUseText{1}{1}{2} und \FDIUseText{1}{1}{4}).}
\FDISetText{3}{3}{2}{\FDIUseText{1}{1}{2}}
\FDISetText{3}{3}{3}{\FDIUseText{1}{1}{4}}
\FDISetText{3}{4}{1}{Lothar selbst\\erhielt das\\Mittelreich.\\Es erstreckte\\sich als langer,\\schmaler\\Korridor von der\\Nordsee bis zum\\Mittelmeer.}
\FDISetText{3}{4}{2}{\FDIUseText{1}{1}{3}}
\FDISetText{3}{5}{1}{Dieses \FDIUseText{1}{1}{3} wurde zur Keimzelle\\der europäischen Integration.}
\FDISetText{3}{6}{1}{Viele Institutionen der\\Europäischen Union sitzen heute entlang\\dieses ehemaligen Korridors\\(Brüssel, Luxemburg, Straßburg, Rom).}

% Seite 4
\FDISetText{4}{1}{1}{So kommt es, dass die Comiciade\\im Technologiezentrum am Europaplatz in\\Aachen ihre Tore für Comicfans aus aller\\Welt öffnet.}
\FDISetText{4}{1}{2}{Sitzen alle\\Sprechblasen an\\der richtigen\\Stelle?}
\FDISetText{4}{1}{3}{Hier geht\\es zur DEUS\\AIX MACHINA-\\Ausstellung!}
\FDISetText{4}{2}{1}{Auch im C.A.R.L. der RWTH Aachen\\laufen hierzu bereits die Vorbereitungen.}
\FDISetText{4}{2}{2}{Der Trend\\geht immer mehr\\in Richtung Data-\\Storytelling!}
\FDISetText{4}{2}{3}{Ist\\Karl der Kleine\\schon einge-\\troffen?}
\FDISetText{4}{3}{1}{Mmmmhh....\\Nicht\\schlecht...}
\FDISetText{4}{4}{1}{Klopf}
\FDISetText{4}{5}{1}{Hallo\\Évariste!}
\FDISetText{4}{6}{1}{Ah,\\da bist\\Du ja,\\Alfred!}

% Seite 5
\FDISetText{5}{1}{1}{Dein Comic über Deine\\Zeit am Lehrstuhl~für~Methodik der\\künstlichen Intelligenz hat mir sehr gut\\gefallen. Was hältst du davon, jetzt den\\ersten FAIRen Comic der Welt zu machen?}
\FDISetText{5}{2}{1}{Den ersten\\FAIRen Comic\\der Welt?}
\FDISetText{5}{3}{1}{Ich meine FAIR im Sinne der\\Forschungsdateninfrastruktur (FDI).}
\FDISetText{5}{3}{2}{Die FAIR-Prinzipien\\sind hierbei:\\Findable, Accessible,\\Interoperable und\\Re-usable.*}
\FDISetText{5}{3}{3}{*Auffindbar, zugänglich, interoperabel und nachnutzbar}
\FDISetText{5}{4}{1}{Das musst Du\\mir näher erklären,\\Évariste.}
\FDISetText{5}{5}{1}{Ganz einfach. Du hast auf der Rückseite\\Deines Comics eine ISBN-Nummer ange-\\geben, damit er im Buchhandel\\leichter zu finden ist.}
\FDISetText{5}{6}{1}{Ja, aus der ISBN-Nummer ist das\\Produkt, der Sprachraum, die Verlags-\\nummer und die Titelnummer\\ablesbar.}
\FDISetText{5}{6}{2}{Produkt\\(Buch)}
\FDISetText{5}{6}{3}{Verlagsnummer\\(Granus Verlag)}
\FDISetText{5}{6}{4}{Titel-\\nummer}
\FDISetText{5}{6}{5}{978-3-9822083-4-3       }
\FDISetText{5}{6}{6}{Ländernummer (Deutschland,\\Österreich und die Schweiz)}
\FDISetText{5}{6}{7}{Prüf-\\ziffer}

% Seite 6
\FDISetText{6}{1}{1}{Man benötigt eine ISBN-Nummer, um in\\das Verzeichnis lieferbarer Bücher (VLB)\\aufgenommen zu werden. Wenn man dort\\nicht gelistet ist, ist es sehr schwer, sein\\Buch in die Läden zu bekommen.}
\FDISetText{6}{2}{1}{Das VLB ist die\\zentrale, unabhängige\\Datenbank und das Recherchetool\\für den deutschsprachigen\\Buchhandel.}
\FDISetText{6}{3}{1}{Interessant. Wir Wissenschaftler\\benutzen anstatt einer ISBN-Nummer\\eine DOI (Digital Object Identifier).}
\FDISetText{6}{4}{1}{Sie ist ein eindeutiger,\\ dauerhafter Identifikator für\\Objekte aller Art. Ein Objekt kann\\mehrere DOIs haben, aber eine\\DOI zeigt immer nur auf\\ein Objekt.}
\FDISetText{6}{4}{2}{DOI-Auflöser}
\FDISetText{6}{4}{3}{Repositorium}
\FDISetText{6}{4}{4}{https://doi.org/10.5281/zenodo.19468332       }
\FDISetText{6}{4}{5}{Präfix}
\FDISetText{6}{4}{6}{Objekt}
\FDISetText{6}{4}{7}{Oder als QR-Code}
\FDISetText{6}{5}{1}{Ein Partner von\\\hspace*{0.75cm} NFDI4ING, der DataCite-Verein,\\vergibt DOIs.*}
\FDISetText{6}{5}{2}{*DataCite ist eine Registrierungsagentur für DOIs}
\FDISetText{6}{6}{1}{DOIs sind bei\\wissenschaftlichen\\Arbeiten der Standard\\für die präzise\\Quellenangabe.}
\FDISetText{6}{6}{2}{Und wie\\genau funk-\\tionieren\\sie?}

% Seite 7
\FDISetText{7}{1}{1}{Sie verweisen auf\\digitale Objekte.}
\FDISetText{7}{1}{2}{Digitales Objekt 1}
\FDISetText{7}{1}{3}{Digitales Objekt 2}
\FDISetText{7}{1}{4}{Digitales Objekt 3}
\FDISetText{7}{1}{5}{Digitales Objekt 4}
\FDISetText{7}{1}{6}{Digitales Objekt 5}
\FDISetText{7}{1}{7}{Digitales Objekt 6}
\FDISetText{7}{2}{1}{Cool. So bleiben wissenschaftliche\\Arbeiten dauerhaft zitierbar, ohne\\dass Links “tot” laufen, oder?}
\FDISetText{7}{3}{1}{Genau so ist es! Jedes digitale Objekt,\\egal ob Artikel, Daten oder Video, erhält\\einen unverwechselbaren Code.}
\FDISetText{7}{4}{1}{Da die DOI nicht\\veränderbar ist, dient sie als\\verlässliche Referenz.}
\FDISetText{7}{5}{1}{Hmmm ...\\Wie kann man\\das am besten\\anschaulich\\machen?}
\FDISetText{7}{6}{1}{Ganz einfach.\\Wir Wissenschaftler\\verwenden hierfür\\gerne die Konserven-\\Metapher.}

% Seite 8
\FDISetText{8}{1}{1}{Bohnen}
\FDISetText{8}{1}{2}{Und so\\funktioniert der\\Vergleich zwischen\\ einer Konserve und\\ digitalem Objekt:\\ Wie auf einer\\Konservendose der\\Inhalt der Dose\\beschrieben wird,\\wird beim Digitalen\\Objekt (DO) der\\Dateninhalt\\beschrieben.}
\FDISetText{8}{2}{1}{Diese Grafik verdeutlicht sehr\\gut wie ein DO aufgebaut ist.}
\FDISetText{8}{2}{2}{Bit-Sequenz}
\FDISetText{8}{2}{3}{Metadaten\\ \& \\Kontext}
\FDISetText{8}{2}{4}{Operations}
\FDISetText{8}{2}{5}{Persistent Identifier}
\FDISetText{8}{3}{1}{Im Zentrum steht die Bit-Sequenz.\\Sie besteht aus Daten in Form von\\Nullen und Einsen.}
\FDISetText{8}{4}{1}{In unserer Konserven-Metapher\\entspricht das dem Inhalt der\\Dose, den Bohnen.}
\FDISetText{8}{5}{1}{Darum herum sind die Metadaten.\\Sie beschreiben den Inhalt in Worten.\\In unserem Fall steht dort\\ \enquote{\FDIUseText{8}{1}{1}}.}
\FDISetText{8}{5}{2}{\FDIUseText{8}{1}{1}}

% Seite 9
\FDISetText{9}{1}{1}{Der nächste Kreis mit den Operations\\erklärt, was man mit den\\Inhalten machen kann.}
\FDISetText{9}{1}{2}{Operations}
\FDISetText{9}{2}{1}{NFDI4ING setzt sich dafür\\ein, mehr als nur eine Download-\\Funktion in den Operations zu\\definieren.}
\FDISetText{9}{3}{1}{Der äußere Kreis ist der Persistent\\Identifier (PID). PID ist der Oberbegriff\\für eine Reihe von Systemen, die digitale\\Objekte auffindbar machen.}
\FDISetText{9}{3}{2}{PID}
\FDISetText{9}{3}{3}{In unserem Fall handelt es sich um eine DOI.}
\FDISetText{9}{4}{1}{Es gibt aber auch andere Systeme, wie\\z.B. URN, Handle oder ORCID.*}
\FDISetText{9}{4}{2}{*URN = Uniform Resource Name, Handle = Handle System; ORCID = für Open Researcher and Contributor ID}
\FDISetText{9}{5}{1}{Jede DOI\\ist also ein PID\\aber nicht jede PID\\ist automatisch\\eine DOI.}
\FDISetText{9}{5}{2}{Uff!\\Du gehst aber\\ganz schön ins\\Eingemachte,\\Évariste!}
\FDISetText{9}{6}{1}{HA HA HA!\\Ja, ich alphabeti-\\siere dich gerade,\\Alfred!}
\FDISetText{9}{6}{2}{Du\\alphabetisierst\\mich?}

% Seite 10
\FDISetText{10}{1}{1}{In der FDI beschreibt der Begriff\\ \enquote{Alphabetisieren} den Prozess des Erwerbs\\grundlegender Fähigkeiten und Kenntnisse,\\die notwendig sind, um Forschungsdaten\\lesen und schreiben zu können.}
\FDISetText{10}{2}{1}{Wenn sie niemand lesen kann, nützt\\auch das beste Archiv nichts und die\\Daten bleiben stumm.}
\FDISetText{10}{2}{2}{Das wusste auch schon\\Karl der Große. Er hat erkannt, dass\\Infrastruktur ohne Bildung nutzlos ist.}
\FDISetText{10}{3}{1}{Die\\Datensprache\\muss man lesen\\lernen.}
\FDISetText{10}{4}{1}{Karl der Große legte großen Wert darauf,\\dass Kleriker und Verwaltungsbeamte\\lesen und schreiben können. Auch, damit\\sie korrekt abschreiben. Fehlerhafte\\Datenübertragung war schon damals\\ein Problem.}
\FDISetText{10}{5}{1}{Sozusagen\\Copy-Paste-Fehler im\\9. Jahrhundert.}
\FDISetText{10}{5}{2}{HA HA HA!}
\FDISetText{10}{5}{3}{Das nenne\\ich doch mal\\eine gelungene\\Interoperabilität von\\Vergangenheit und\\Gegenwart!}
\FDISetText{10}{6}{1}{Was gibt\\es denn hier zu\\lachen?}

% Seite 11
\FDISetText{11}{1}{1}{Hallo Tobias. Wir sind gerade dabei,\\einen FAIREN Comic zu gestalten, und\\das ist gerade sehr unterhaltsam.}
\FDISetText{11}{2}{1}{Das ist ja spannend. Ich habe gestern mein\\neuestes Paper bei ing.grid eingereicht.\\Das könntet ihr doch mit diesem Comic\\auch machen.}
\FDISetText{11}{2}{2}{ing.grid?\\Was ist das\\denn?}
\FDISetText{11}{3}{1}{ing.grid ist ein\\Diamond-Open-Access-\\Journal der NFDI4ING.}
\FDISetText{11}{3}{2}{Dort\\gibt es\\auch einen\\Preprint-\\Server.}
\FDISetText{11}{4}{1}{Das ist eine tolle Idee, Tobias.\\ Wenn wir dort einreichen, wäre der Comic\\schon auf dem Preprintserver sichtbar\\und für alle zugägnlich.}
\FDISetText{11}{5}{1}{Genau! Ein Community-Review wäre\\dort der erste Schritt. Das heißt, der\\Comic könnte dann bereits kommentiert\\und bewertet werden.}
\FDISetText{11}{6}{1}{Super!\\Wo soll ich meinen\\Vollziehungsstrich\\hinsetzen?}

% Seite 12
\FDISetText{12}{1}{1}{Einen Moment noch!\\Zuerst müssen wir den Comic in die\\Zielsprache übersetzen. Dafür gibt es\\auch schon eine Handreichung.}
\FDISetText{12}{2}{1}{Diese Handreichung findet\\man unter diesem HANDLE.}
\FDISetText{12}{2}{2}{http://hdl.handle.net/2268.2/12918}
\FDISetText{12}{2}{3}{Als Zitation sieht das so aus: \cite{simons_traduction_2021}}
\FDISetText{12}{3}{1}{Okay.\\Dann laden\\wir ihn als PDF\\hoch.*}
\FDISetText{12}{3}{2}{* PDF = Portable Document Format}
\FDISetText{12}{4}{1}{Uuuund...\\Los gehts!}
\FDISetText{12}{4}{2}{KLICK}
\FDISetText{12}{5}{1}{Ihre Datei wurde erfolgreich hochgeladen.}
\FDISetText{12}{5}{2}{PLING!}
\FDISetText{12}{6}{1}{Jetzt wird es\\spannend!}
\FDISetText{12}{6}{2}{Ende des\\ersten Teils}

